\section*{自同构}

\begin{enumerate}
\def\labelenumi{\arabic{enumi})}
\setcounter{enumi}{15}
\tightlist
\item
  \(\mathrm{Aut}(\mathfrak{g})\) 中由诸 \(e^{\mathrm{ad}x}\)（\(x\) 幂零）生成的子群是初等自同构群 \(\mathrm{Aut}_e(\mathfrak{g})\)（\(\mathfrak{g}\) 的）；它是 \(\mathrm{Aut}(\mathfrak{g})\) 的正规子群；它等于自身的换位子群。
\end{enumerate}

若 \(\bar{k}\) 是 \(k\) 的代数闭包，则群 \(\operatorname{Aut}(\mathfrak{g})\) 自然地嵌入 \(\operatorname{Aut}(\mathfrak{g} \otimes_k \bar{k})\)。置

\[
\operatorname{Aut} _ {0} (\mathfrak {g}) = \operatorname{Aut} (\mathfrak {g}) \cap \operatorname{Aut} _ {e} (\mathfrak {g} \otimes_ {k} \bar {k});
\]

这是 \(\operatorname{Aut}(\mathfrak{g})\) 的正规子群，与 \(\bar{k}\) 的选取无关。我们有

\[
\operatorname{Aut} _ {e} (\mathfrak {g}) \subset \operatorname{Aut} _ {0} (\mathfrak {g}) \subset \operatorname{Aut} (\mathfrak {g}).
\]

\(\mathrm{Aut}_0(\mathfrak{g})\) 的换位子群是 \(\mathrm{Aut}_e(\mathfrak{g})\)。在 Zariski 拓扑下，\(\mathrm{Aut}(\mathfrak{g})\) 与 \(\mathrm{Aut}_0(\mathfrak{g})\) 在 \(\mathrm{End}_k(\mathfrak{g})\) 中闭，\(\mathrm{Aut}_0(\mathfrak{g})\) 是 \(\mathrm{Aut}(\mathfrak{g})\) 的单位元连通分支，且 \(\mathrm{Aut}_e(\mathfrak{g})\) 在 \(\mathrm{Aut}_0(\mathfrak{g})\) 中稠密。

令 B 为 R 的基，\(\operatorname{Aut}(R,B)\) 为使 B 稳定的 R 的自同构群。则 \(\operatorname{Aut}(\mathfrak{g})\) 是一个同构于 \(\operatorname{Aut}(R,B)\) 的子群与 \(\operatorname{Aut}_{0}(\mathfrak{g})\) 的半直积；特别地，\(\operatorname{Aut}(\mathfrak{g})/\operatorname{Aut}_{0}(\mathfrak{g})\) 同构于 \(\operatorname{Aut}(R,B)\)，而后者本身又同构于 g 的邓金图的一个自同构群。

\begin{enumerate}
\def\labelenumi{\arabic{enumi})}
\setcounter{enumi}{16}
\tightlist
\item
  g 的标架是指一个三元组 \((\mathfrak{h}^{\prime},\mathrm{B},(X_{\alpha})_{\alpha\in\mathrm{B}})\)，其中 \(h^{\prime}\) 是 g 的分裂嘉当子代数，B 是 \(\mathrm{R}(\mathfrak{g},\mathfrak{h}^{\prime})\) 的基，且对所有 \(\alpha\in B\)，\(X_{\alpha}\) 是 \(g^{\alpha}\) 的非零元素。群 \(\operatorname{Aut}_{0}(\mathfrak{g})\) 单迁移地作用于 g 的标架所成的集合上。
\end{enumerate}

群 \(\mathrm{Aut}_e(\mathfrak{g})\) 迁移地作用于偶对 \((\mathfrak{k},\mathfrak{b})\) 所成的集合上，其中 \(\mathfrak{k}\) 是 \(\mathfrak{g}\) 的分裂嘉当子代数，\(\mathfrak{b}\) 是 \((\mathfrak{g},\mathfrak{k})\) 的 Borel 子代数。

\begin{enumerate}
\def\labelenumi{\arabic{enumi})}
\setcounter{enumi}{17}
\tightlist
\item
  以 \(\operatorname{Aut}(\mathfrak{g},\mathfrak{h})\) 表示由 \(s\in \operatorname{Aut}(\mathfrak{g})\)（满足 \(s(\mathfrak{h}) = \mathfrak{h}\)）组成的集合。置
\end{enumerate}

\[
\operatorname{Aut} (\mathfrak {g}, \mathfrak {h}) = \operatorname{Aut} _ {e} (\mathfrak {g}) \cap \operatorname{Aut} (\mathfrak {g}, \mathfrak {h}), \quad \operatorname{Aut} _ {0} (\mathfrak {g}, \mathfrak {h}) = \operatorname{Aut} _ {0} (\mathfrak {g}) \cap \operatorname{Aut} (\mathfrak {g}, \mathfrak {h}).
\]

若 \(s \in \text{Aut}(\mathfrak{g}, \mathfrak{h})\)，则 \(s | \mathfrak{h}\) 的反步映射是 R 的自同构群 \(\text{A(R)}\) 的一个元素；以 \(\varepsilon(s)\) 记这个元素；映射 \(\varepsilon\) 是从 \(\text{Aut}(\mathfrak{g}, \mathfrak{h})\) 到 \(\text{A(R)}\) 的同态。我们有 \(\text{Aut}_{0}(\mathfrak{g}) = \text{Aut}_{e}(\mathfrak{g})\). Ker \(\varepsilon\)，且

\[
\varepsilon (\operatorname{Aut} _ {0} (\mathfrak {g}, \mathfrak {h})) = \varepsilon (\operatorname{Aut} _ {e} (\mathfrak {g}, \mathfrak {h})) = W.
\]

令 \(\mathrm{T}_{\mathrm{P}} = \mathrm{Hom}(\mathrm{P}(\mathrm{R}), k^{*})\)，\(\mathrm{T}_{\mathrm{Q}} = \mathrm{Hom}(\mathrm{Q}(\mathrm{R}), k^{*})\)。\(\mathrm{Q}(\mathrm{R})\) 到 \(\mathrm{P}(\mathrm{R})\) 的单射定义了一个从 \(\mathrm{T}_{\mathrm{P}}\) 到 \(\mathrm{T}_{\mathrm{Q}}\) 的同态；以 \(\mathrm{Im}(\mathrm{T}_{\mathrm{P}})\) 表示其像。若 \(t \in \mathrm{T}_{\mathrm{Q}}\)，令 \(f(t)\) 为 \(\mathfrak{g}\) 的满足下述条件的自同态：对所有 \(\alpha \in \mathrm{R} \cup \{0\}\)，\(f(t)|\mathfrak{g}^{\alpha}\) 是比为 \(t(\alpha)\) 的位似；我们有 \(f(t)\in \mathrm{Aut}_{0}(\mathfrak{g},\mathfrak{h})\)，且 \(f\) 是从 \(\mathrm{T}_{\mathrm{Q}}\) 到 \(\mathrm{Aut}_0(\mathfrak{g},\mathfrak{h})\) 的单同态。序列

\[
\{1 \} \longrightarrow \mathrm {T_ {Q}} \stackrel {{f}} {{\longrightarrow}} \operatorname{Aut} (\mathfrak {g}, \mathfrak {h}) \stackrel {{\varepsilon}} {{\longrightarrow}} \mathrm{A(R)} \longrightarrow \{1 \}
\]

与

\[
\{1 \} \longrightarrow \mathrm{T} _ {\mathrm{Q}} \stackrel {f} {\longrightarrow} \mathrm{Aut} _ {0} (\mathfrak {g}, \mathfrak {h}) \stackrel {\varepsilon} {\longrightarrow} \mathrm{W} \longrightarrow \{1 \}
\]

是正合的。我们有 \(f(\mathrm{Im}(\mathrm{T}_{\mathrm{P}})) \subset \mathrm{Aut}_{e}(\mathfrak{g}, \mathfrak{h})\)；f 通过过渡到商定义出一个满射\(^{23}\)同态 \(\mathrm{T}_{\mathrm{Q}}/\mathrm{Im}(\mathrm{T}_{\mathrm{P}}) \to \mathrm{Aut}_{0}(\mathfrak{g})/\mathrm{Aut}_{e}(\mathfrak{g})\)。在 Zariski 拓扑下，\(f(\mathrm{T}_{\mathrm{Q}})\) 在 \(\mathrm{Aut}(\mathfrak{g})\) 中闭，且 \(f(\mathrm{Im}(\mathrm{T}_{\mathrm{P}}))\) 在 \(f(\mathrm{T}_{\mathrm{Q}})\) 中稠密。

